\clearpage
\SourceLeaf{175}%
\pdfbookmark[0]{Göttingische gelehrta Anzeigen. 1825 April 11.}{notice1825april11}
\begin{center}
\rule{.82\linewidth}{.35pt}\par
.Göttingische gelehrta Anzeigen.\quad 1825 April 11.\par
\rule{.82\linewidth}{.35pt}
\end{center}
Am 5. April überreichte Hr. Hofr. \textsc{Gauss} der Königl. Societät eine Vorlesung, überschrieben:
\begin{center}\itshape
Theoria Residuorum Biquadraticorum, Commentatio prima.
\end{center}

Die Theorie der quadratischen Reste bildet bekanntlich einen der interessantesten Theile der Höhern Arithmetik, welchen man jetzt nach vielfach wiederholten Untersuchungen als vollendet und abgeschlossen betrachten kann: die Geschichte desselben betreffende Nachrichten findet man in diesen Blättern 1808 Mai 12 und Sept. 19, und 1817 März 10. An letzterm Orte sind auch bereits einige vorläufige Nachrichten über die Nachforschungen mitgetheilt, welche der Verfasser der vorliegenden Abhandlung seit dem Jahre 1805 über die verwandte, eben so fruchtbare und interessante, aber sehr viel schwierigere Theorie der cubischen und biquadratischen Reste angestellt hatte. Obgleich schon damals im Besitz der wesentlichen Momente dieser Theorie, ist er doch bisher durch andere Arbeiten abgehalten, öffentlich etwas davon bekannt zu machen, und erst jetzt ist es ihm möglich geworden, sich mit der Ausarbeitung eines Theils dieser Untersuchungen zu beschäftigen. Der Anfang ist jetzt mit der Theorie der biquadratischen Reste gemacht, die der Theorie der quadratischen Reste näher verwandt ist, als die der cubischen. Inzwischen ist die gegenwärtige Abhandlung
\SourceLeaf{176}%
noch keinesweges dazu bestimmt, den überaus reichhaltigen Gegenstand zu erschöpfen. Die Entwickelung der \emph{allgemeinen} Theorie, welche eine ganz eigenthümliche Erweiterung des Feldes der höhern Arithmetik erfordert, bleibt vielmehr der künftigen Fortsetzung vorbehalten, während in diese erste Abhandlung diejenigen Untersuchungen aufgenommen sind, welche sich ohne eine solche Erweiterung vollständig darstellen liessen. Von den Resultaten kann in dieser Anzeige nur ein Theil ausgehoben werden.

Eine ganze Zahl $a$ heisst biquadratischer Rest der ganzen Zahl $p$, wenn es Zahlen der Form $x^4-a$ gibt, die durch $p$ theilbar sind; biquadratischer Nichtrest hingegen, wenn keine Zahlen jener Form durch $p$ theilbar sein können. Offenbar sind alle biquadratischen Reste von $p$ zugleich quadratische Reste derselben Zahl, und also alle quadratischen Nichtreste auch biquadratische Nichtreste: allein nicht alle quadratischen Reste sind zugleich biquadratische Reste. Es ist zureichend, die Untersuchungen auf den Fall einzuschränken, wo $p$ eine Primzahl von der Form $4n+1$, und $a$ nicht durch $p$ theilbar ist, da alle anderen Fälle sich leicht auf diesen zurückführen lassen.

Die Untersuchungen über diesen Gegenstand zerfallen in zwei Abtheilungen, je nachdem $p$ oder $a$ als gegeben angesehen wird. Die erstere ist von viel geringerer Schwierigkeit als die zweite, und verglichen mit letzterer als ganz elementarisch zu betrachten. Alles Wesentliche, was darüber zu sagen ist, enthält die Abhandlung vollständig.

Aus der zweiten Abtheilung hingegen sind hier nur erst einige specielle Fälle abgehandelt, die sich ohne zu grosse Zurüstungen abmachen liessen, und als Vorbereitungen zu der künftig zu gebenden allgemeinen Theorie dienen können. Dies sind diejenigen, wo $a=-1$, und $a=\pm2$ gesetzt wird. Der erstere Fall hat gar keine Schwierigkeit: es war auch schon in dem Werke, \emph{Disquisitiones Arithmeticae}, gezeigt, dass $-1$ ein biquadratischer Rest von $p$ ist, so oft $p$ die Form $8n+1$ hat, hingegen ein blos quadratischer Rest und biquadratischer Nichtrest von $p$, wenn $p$ von der Form $8n+5$ wird. Ganz anders verhält es sich mit dem Fall $a=\pm2$. Es ist zwar längst bekannt, dass $+2$ und $-2$ von $p$ quadratische und also auch biquadratische Nichtreste sind, wenn $p$ die Form $8n+5$ hat, und wenigstens quadratische Reste, wenn $p$ von der Form $8n+1$ ist, wie auch dass bei dieser Form von $p$ entweder $+2$ und $-2$ zugleich biquadratische Reste, oder zugleich biquadratische Nichtreste werden: 
\SourceLeaf{177}%
allein die Unterscheidung, welcher dieser beiden Fälle eintrete, ist eine Untersuchung von viel höherer Art, und es werden dazu in der Abhandlung zwei verschiedene Criterien entwickelt.

Das erste Criterium hängt mit der Zerlegung der Zahl $p$ in ein einfaches und ein doppeltes Quadrat zusammen, die bekanntlich (da, wie schon bemerkt ist, angenommen wird, dass $p$ eine Primzahl sei) immer möglich und nur auf Eine Art möglich ist. Setzt man $p=gg+2hh$, so wird $\pm2$ ein biquadratischer Rest von $p$, wenn $g$ von der Form $8n+1$ oder $8n+7$, ein biquadratischer Nichtrest hingegen, wenn $g$ von der Form $8n+3$ oder $8n+5$ ist.

Das zweite Criterium hängt zusammen mit der Zerlegung der Zahl $p$ in zwei Quadrate, die bekanntlich auch immer möglich und nur auf Eine Art möglich ist. Setzt man $p=ee+ff$, und nimmt an, dass $ee$ das ungerade, $ff$ das gerade Quadrat bedeutet, so bringt schon die vorausgesetzte Form von $p=8n+1$ mit sich, dass auch $\frac12f$ eine gerade Zahl wird, also $f$ entweder von der Form $8m$ oder von der Form $8m+4$: im ersten Fall nun wird $\pm2$ biquadratischer Rest, im andern biquadratischer Nichtrest von $p$ sein.

Wir deuten hier nur die Bemerkung an, wozu die höhere Arithmetik so oft Gelegenheit gibt, dass nicht so wohl die Schönheit und Einfachheit der Theoreme selbst, als die Schwierigkeit ihrer Begründung sie vorzüglich merkwürdig macht. Sobald man einmal veranlasst ist, das Dasein eines Zusammenhanges zwischen dem Verhalten der Zahl $\pm2$ und den beiden angeführten Zerlegungen der Zahl $p$ zu vermuthen, ist es äusserst leicht, diesen Zusammenhang durch Induction wirklich zu entdecken. Allein schon bei dem ersten Criterium ist der Beweis dafür nicht ganz leicht zu führen, viel tiefer versteckt liegt er aber bei dem zweiten, wo er mit anderweitigen subtilen Hülfsuntersuchungen innigst verkettet ist, die ihrerseits wieder zu einer merkwürdigen Erweiterung der Theorie der Kreistheilung führen. Diese wunderbare Verkettung der Wahrheiten ist es vorzüglich. was, wie man schon oft bemerkt hat, der höhern Arithmetik einen so eigenthümlichen Reiz gibt. Diese Begründungen selbst vertragen übrigens natürlich hier keinen Auszug, und müssen in der Abhandlung selbst nachgesehen werden. Allein ein paar andere neue arithmetische Theoreme, welche gleichfalls mit der Begründung des zweiten Criterium innigst verbunden sind, verdienen wohl, ihrer Einfachheit wegen, hier noch besonders herausgehoben zu werden.

Wenn $p$ eine Primzahl von der Form $4k+1$ ist, und $=ee+ff$ 
\SourceLeaf{178}%
gesetzt wird, so dass $ee$ das ungerade, $ff$ das gerade Quadrat bedeutet; wenn man ferner
\[
1\pdot2\pdot3\mathinner{.\,.\,.\,.\,.\,.}\ k=q
\]
\[
(k+1)(k+2)(k+3)\dotsfive2k=r
\]
setzt, so wird allemal $\pm e$ der kleinste Rest sein, welcher hervorgeht, indem man $\frac r{2q}$ mit $p$ dividirt, und $\pm f$ der kleinste Rest, welchen man aus der Division von $\frac12rr$ mit $p$ erhält (kleinsten Rest immer so verstanden, dass er zwischen den Grenzen $-\frac12p$ und $+\frac12p$ genommen wird). Die Zahl $\frac r{2q}$, welche für $p=5$ den Werth $1$ erhält, kann man für grössere Werthe von $p$ auch in folgende Form setzen
\[
\frac{6\pdot10\pdot14\pdot18\dotsfive(p-3)}{2\pdot3\pdot4\pdot5\dotsfive k}
\]
Es ist sehr merkwürdig, dass so die Zerlegung der Zahl $p$ in zwei Quadrate ganz auf directem Wege erhalten werden kann: aber fast noch merkwürdiger ist ein dabei Statt findender Nebenumstand. Allemal nemlich findet man durch dieses Verfahren die Wurzel des ungeraden Quadrates, $e$, mit positivem Zeichen, wenn $e$, positiv genommen, von der Form $4m+1$ ist, und mit negativem, wenn $e$ positiv genommen von der Form $4m+3$ ist. Hingegen hat für das Zeichen, mit welchem die Wurzel des geraden Quadrats, $f$, aus jener Operation hervorgeht, noch durchaus keine allgemeine Regel aufgefunden werden können, weder a priori, noch auf dem Wege der Induction, und der Verfasser empfiehlt daher, am Schlusse der Abhandlung, diesen Gegenstand den Freunden der höhern Arithmetik zu weiterer Nachforschung, überzeugt, dass mit dem Gelingen derselben sich zugleich eine ergiebige Quelle neuer Erweiterungen dieses schönen Theils der Mathematik eröffnen werde.
\BookRule{.28\linewidth}
